\section{Echo dan Print: Menampilkan Output}

\textbf{Echo} dan \textbf{print} adalah dua cara utama untuk menampilkan output dari PHP ke browser. Meskipun fungsinya serupa, terdapat perbedaan teknis di antara keduanya. \code{echo} dapat menerima beberapa argumen yang dipisah koma dan tidak mengembalikan nilai (void), sedangkan \code{print} hanya menerima satu argumen dan selalu mengembalikan nilai \code{1}. Dalam praktik sehari-hari, \code{echo} lebih sering digunakan karena sedikit lebih cepat dan lebih fleksibel \cite{phpManual}.

Kedua perintah dapat dipakai dengan atau tanpa tanda kurung: \code{echo "Halo";} sama dengan \code{echo("Halo");}. Untuk menampilkan string yang mengandung variabel, cara paling mudah adalah menggunakan kutip ganda (\textit{double quotes}) yang menginterpretasikan variabel secara langsung, atau menggunakan operator konkatenasi (titik). String dengan kutip tunggal (\textit{single quotes}) tidak menginterpretasikan variabel dan hanya mengenali escape \code{\textbackslash\textbackslash} dan \code{\textbackslash'} \cite{w3schoolsPhp}.

\begin{table}[ht]
\centering
\begin{tabular}{|P{3cm}|P{4.5cm}|P{4.5cm}|}
\hline
\textbf{Aspek} & \textbf{echo} & \textbf{print} \\
\hline
Jumlah argumen & Banyak (dipisah koma) & Satu saja \\
\hline
Nilai kembalian & Tidak ada (void) & Selalu \code{1} \\
\hline
Kecepatan & Sedikit lebih cepat & Sedikit lebih lambat \\
\hline
Penggunaan di ekspresi & Tidak bisa & Bisa (karena mengembalikan nilai) \\
\hline
Tanda kurung & Opsional & Opsional \\
\hline
\end{tabular}
\caption{Perbandingan echo dan print di PHP}
\end{table}

\begin{webcode}[language=PHP, caption={Berbagai cara menggunakan echo dan print}]
<?php
$nama = "Dewi";
$umur = 22;

// echo: beberapa argumen
echo "Nama: ", $nama, ", Umur: ", $umur, "\n";

// echo: konkatenasi
echo "Halo, " . $nama . "!\n";

// echo: interpolasi variabel (kutip ganda)
echo "Hai $nama, umur Anda $umur tahun.\n";

// print: satu argumen, return 1
$hasil = print "Ini print.\n";
echo "Return print: $hasil\n";  // 1

// echo shortcut di dalam HTML
?>
<p>Nama: <?= $nama ?></p>  <!-- sama dengan <?php echo $nama; ?> -->
\end{webcode}

Untuk keperluan debugging, PHP menyediakan fungsi \code{var\_dump()} dan \code{print\_r()} yang menampilkan informasi lebih detail tentang variabel. Fungsi \code{var\_dump()} menampilkan tipe dan nilai variabel, sehingga sangat berguna saat mencari bug terkait tipe data. Fungsi \code{print\_r()} menampilkan nilai variabel dalam format yang lebih mudah dibaca, terutama untuk array dan objek. Kedua fungsi ini sebaiknya hanya digunakan saat pengembangan, bukan di lingkungan produksi \cite{phpRightWay}.

\begin{catatan}
\textbf{Keamanan Output:} Saat menampilkan data yang berasal dari input pengguna (misalnya dari form atau URL), \textbf{selalu} gunakan \code{htmlspecialchars(\$data, ENT\_QUOTES, 'UTF-8')} untuk mencegah serangan \textit{Cross-Site Scripting} (XSS). Tanpa sanitasi, penyerang dapat menyuntikkan kode JavaScript berbahaya melalui input form \cite{owaspXss}.
\end{catatan}

